% !Mode:: "TeX:UTF-8" 
\begin{conclusions}

本文围绕基于 FH-MFSK 的声学释放器及测距技术开展了系统研究，针对浅海多径、噪声和时变环境下声学释放器通信可靠性不足、测距稳定性受限等问题，从浅海信道分析、通信算法设计、测距模型构建以及软硬件系统实现四个层面进行了研究与归纳。本文完成的主要工作可总结如下：

（1）分析了浅海水声信道的多径扩展、噪声干扰、多普勒效应和频率选择性衰落特征，明确了浅海复杂环境对 FH-MFSK 体制在跳速、频点间隔、同步方式和判决门限等方面的约束关系。

（2）设计了面向声学释放器控制链路的 FH-MFSK 通信算法，建立了包含前导、同步、地址、命令和校验字段的帧结构，研究了基于能量门限与相关统计量的同步捕获方法，以及适用于低复杂度嵌入式平台的非相干检测流程。

（3）建立了基于往返时延的测距模型，给出了首达径检测、多径抑制、声速修正和异常值剔除等关键方法，形成了适合浅海声学释放器工程应用的测距处理流程。

（4）完成了声学释放器总体方案、水下机与甲板单元硬件划分、嵌入式软件流程和测试方案设计，为后续开展实验室联调、水池试验和海上外场验证提供了系统实现基础。

受论文撰写阶段实验数据整理进度限制，部分性能结果仍需结合后续实测数据进一步量化。下一步工作可从以下几个方向展开：一是结合真实海试数据对同步门限、跳频序列和频点配置进行优化；二是引入更精细的声速修正模型，提高中远距离测距精度；三是围绕低功耗、小型化和高可靠执行机构开展进一步工程化设计；四是推进整机产品化验证，形成面向实际海洋作业的应用规范。

\end{conclusions}
